\documentclass[10pt,a4paper]{amsart}
\usepackage{amsmath,amssymb,amsthm,graphicx}
% Retained minimal preamble for the Stacks Project chapter varieties.tex.
% Verbatim excerpt of the chapter preamble: Xy-pic setup (native preamble lines 14-18), the native theorem
% declarations with the mathematical macros (native preamble lines 154-217), and the project own no-op
% definitions of the reference/slogan/history display environments (native wrapper.tex lines 1-3) that the
% counted statement uses for its slogan. Copied mechanically; nothing added, removed or redefined.
% For commutative diagrams we use Xy-pic
\usepackage[all]{xy}

% We use 2cell for 2-commutative diagrams.
\xyoption{2cell}
% Theorem environments.
%
\theoremstyle{plain}
\newtheorem{theorem}[subsection]{Theorem}
\newtheorem{proposition}[subsection]{Proposition}
\newtheorem{lemma}[subsection]{Lemma}

\theoremstyle{definition}
\newtheorem{definition}[subsection]{Definition}
\newtheorem{example}[subsection]{Example}
\newtheorem{exercise}[subsection]{Exercise}
\newtheorem{situation}[subsection]{Situation}

\theoremstyle{remark}
\newtheorem{remark}[subsection]{Remark}
\newtheorem{remarks}[subsection]{Remarks}

\numberwithin{equation}{subsection}

% Macros
%
\def\lim{\mathop{\mathrm{lim}}\nolimits}
\def\colim{\mathop{\mathrm{colim}}\nolimits}
\def\Spec{\mathop{\mathrm{Spec}}}
\def\Hom{\mathop{\mathrm{Hom}}\nolimits}
\def\Ext{\mathop{\mathrm{Ext}}\nolimits}
\def\SheafHom{\mathop{\mathcal{H}\!\mathit{om}}\nolimits}
\def\SheafExt{\mathop{\mathcal{E}\!\mathit{xt}}\nolimits}
\def\Sch{\mathit{Sch}}
\def\Mor{\mathop{\mathrm{Mor}}\nolimits}
\def\Ob{\mathop{\mathrm{Ob}}\nolimits}
\def\Sh{\mathop{\mathit{Sh}}\nolimits}
\def\NL{\mathop{N\!L}\nolimits}
\def\CH{\mathop{\mathrm{CH}}\nolimits}
\def\proetale{{pro\text{-}\acute{e}tale}}
\def\etale{{\acute{e}tale}}
\def\QCoh{\mathit{QCoh}}
\def\Ker{\mathop{\mathrm{Ker}}}
\def\Im{\mathop{\mathrm{Im}}}
\def\Coker{\mathop{\mathrm{Coker}}}
\def\Coim{\mathop{\mathrm{Coim}}}

% Boxtimes
%
\DeclareMathSymbol{\boxtimes}{\mathbin}{AMSa}{"02}

%
% Macros for moduli stacks/spaces
%
\def\QCohstack{\mathcal{QC}\!\mathit{oh}}
\def\Cohstack{\mathcal{C}\!\mathit{oh}}
\def\Spacesstack{\mathcal{S}\!\mathit{paces}}
\def\Quotfunctor{\mathrm{Quot}}
\def\Hilbfunctor{\mathrm{Hilb}}
\def\Curvesstack{\mathcal{C}\!\mathit{urves}}
\def\Polarizedstack{\mathcal{P}\!\mathit{olarized}}
\def\Complexesstack{\mathcal{C}\!\mathit{omplexes}}
% \Pic is the operator that assigns to X its picard group, usage \Pic(X)
% \Picardstack_{X/B} denotes the Picard stack of X over B
% \Picardfunctor_{X/B} denotes the Picard functor of X over B
\def\Pic{\mathop{\mathrm{Pic}}\nolimits}
\def\Picardstack{\mathcal{P}\!\mathit{ic}}
\def\Picardfunctor{\mathrm{Pic}}
\def\Deformationcategory{\mathcal{D}\!\mathit{ef}}
\newenvironment{reference}{}{}
\newenvironment{slogan}{}{}
\newenvironment{history}{}{}
\usepackage{cleveref}
\crefname{originallemma}{original lemma}{original lemmas}
\Crefname{originallemma}{Original lemma}{Original lemmas}
\makeatletter
\newlabel{algebra-lemma-p-ring-map}{{0BRA}{source}{}{sourcepointer-1}{}}
\newlabel{algebra-lemma-p-ring-map@cref}{{[originallemma][0][]0BRA}{[1][0][]source}{}{}{}}
\makeatother
\crefname{originallemma}{original lemma}{original lemmas}
\Crefname{originallemma}{Original lemma}{Original lemmas}
\makeatletter
\newlabel{algebra-lemma-minimal-prime-reduced-ring}{{00EU}{source}{}{sourcepointer-2}{}}
\newlabel{algebra-lemma-minimal-prime-reduced-ring@cref}{{[originallemma][0][]00EU}{[1][0][]source}{}{}{}}
\makeatother
\crefname{originallemma}{original lemma}{original lemmas}
\Crefname{originallemma}{Original lemma}{Original lemmas}
\makeatletter
\newlabel{algebra-lemma-make-separable}{{030R}{source}{}{sourcepointer-3}{}}
\newlabel{algebra-lemma-make-separable@cref}{{[originallemma][0][]030R}{[1][0][]source}{}{}{}}
\makeatother
\crefname{originallemma}{original lemma}{original lemmas}
\Crefname{originallemma}{Original lemma}{Original lemmas}
\makeatletter
\newlabel{algebra-lemma-characterize-separable-field-extensions}{{030W}{source}{}{sourcepointer-4}{}}
\newlabel{algebra-lemma-characterize-separable-field-extensions@cref}{{[originallemma][0][]030W}{[1][0][]source}{}{}{}}
\makeatother
\crefname{originallemma}{original lemma}{original lemmas}
\Crefname{originallemma}{Original lemma}{Original lemmas}
\makeatletter
\newlabel{lemma-generic-points-geometrically-reduced}{{04KS}{source}{}{sourcepointer-5}{}}
\newlabel{lemma-generic-points-geometrically-reduced@cref}{{[originallemma][0][]04KS}{[1][0][]source}{}{}{}}
\makeatother
\title{A finite purely inseparable extension makes the reduction of a finite type scheme geometrically reduced}
\author{The Stacks Project Authors}
\date{stacks-project:varieties.tex:local-snapshot-20261011; GNU FDL 1.2 or later}
\begin{document}
\maketitle
\noindent\textit{Editorial note.} All delivered mathematical text is the author's own, copied from the pinned Stacks Project chapter file varieties.tex; no statement or proof step was rewritten, reordered or removed.\par\medskip
\noindent\textit{Editorial note.} The preamble is the chapter's own retained minimal preamble (sources/stacks-varieties/retained-preamble.tex, lines 1-77), a verbatim excerpt of the native Xy-pic setup, of the native theorem declarations with mathematical macros, and of the project's own no-op definitions of the reference, slogan and history display environments; no macro was invented or redefined.\par\medskip
\noindent\textit{Editorial note.} The same-chapter reference is located at its \textbackslash{}label line inside the pinned chapter file; the four cross-chapter pointers are located in the project's own tag list staged in the snapshot (sources/stacks-chow/tags.txt), at the verbatim lines carrying their canonical tags.\par\medskip
\noindent\textit{Editorial note.} Pre-open checks run for this card: (1) same-name intersection filter across staged chapters (no hit); (2) cross-lane scan over every card in focused-seeds and the preparation worklist via agent-authoring/stk/card\_plan.py --scan (no card claims this label).\par\medskip
\noindent\textit{Allowed standard prerequisites (editorial).}\begin{itemize}
\item A purely inseparable extension does not change the underlying topological space of the reduction. Algebra, Lemma labelled algebra-lemma-p-ring-map (canonical Stacks tag 0BRA); invoked at proof lines 651-653. Supplies that (X\_\{k'\})\_\{red\} -> X is a homeomorphism, so the points over the generic points of the components are again the generic points.
\item Local rings at the generic points of a reduced scheme are fields. Algebra, Lemma labelled algebra-lemma-minimal-prime-reduced-ring (canonical Stacks tag 00EU); invoked at proof lines 656-657. Supplies that K\_i = O\_\{X,x\_i\} is a field, and (with local finiteness of type) that K\_i/k is finitely generated.
\item Making finitely generated field extensions separable by a finite purely inseparable base change. Algebra, Lemma labelled algebra-lemma-make-separable (canonical Stacks tag 030R); invoked at proof lines 661-664. Supplies the single finite purely inseparable k'/k that makes all (K\_i tensor\_k k')\_\{red\} separable over k', hence geometrically reduced by algebra-lemma-characterize-separable-field-extensions (030W, line 667).
\item Geometric reducedness is detected at the generic points of the irreducible components. Same chapter, Lemma labelled lemma-generic-points-geometrically-reduced (source line 555, Stacks tag 04KS), part (3); invoked at proof lines 668-670. Supplies the final step from the local rings at the generic points to geometric reducedness of the whole scheme.
\end{itemize}\medskip
\ifdefined\hypertarget\hypertarget{sourcepointer-1}{}\fi
\noindent\textit{Source reference.} Original lemma 0BRA: sources/stacks-chow/tags.txt, source lines 14763--14763. Cross-chapter standard prerequisite invoked once at proof line 653 for the homeomorphism (X\_\{k'\})\_\{red\} -> X; the generic-point bookkeeping on top of it is retained in the counted span.\par\medskip
\ifdefined\hypertarget\hypertarget{sourcepointer-2}{}\fi
\noindent\textit{Source reference.} Original lemma 00EU: sources/stacks-chow/tags.txt, source lines 548--548. Cross-chapter standard prerequisite invoked once at proof lines 656-657 to see that the local rings at generic points of a reduced scheme are fields.\par\medskip
\ifdefined\hypertarget\hypertarget{sourcepointer-3}{}\fi
\noindent\textit{Source reference.} Original lemma 030R: sources/stacks-chow/tags.txt, source lines 3958--3958. Cross-chapter standard prerequisite invoked once at proof lines 661-664 for the single finite purely inseparable extension making all the residue fields separable; the rest of the proof is retained in the counted span.\par\medskip
\ifdefined\hypertarget\hypertarget{sourcepointer-4}{}\fi
\noindent\textit{Source reference.} Original lemma 030W: sources/stacks-chow/tags.txt, source lines 3963--3963. Cross-chapter standard prerequisite invoked once at proof line 667 for separable implies geometrically reduced over the base field.\par\medskip
\ifdefined\hypertarget\hypertarget{sourcepointer-5}{}\fi
\noindent\textit{Source reference.} Original lemma 04KS: sources/stacks-varieties/pinned-original-source.tex, source lines 555--555. Same-chapter pointer invoked at proof lines 668-670, part (3), as the criterion that upgrades the local checks at generic points to geometric reducedness of the scheme; it is the closing step only.\par\medskip
\begin{lemma}
\label{lemma-finite-extension-geometrically-reduced}
Let $k$ be a field.
Let $X \to \Spec(k)$ be locally of finite type.
Assume $X$ has finitely many irreducible components.
Then there exists a finite purely inseparable extension $k'/k$
such that $(X_{k'})_{red}$ is geometrically reduced over $k'$.
\end{lemma}

\begin{proof}
To prove this lemma we may replace $X$ by its reduction $X_{red}$.
Hence we may assume that $X$ is reduced and locally of finite type
over $k$.
Let $x_1, \ldots, x_n \in X$ be the generic points of the irreducible
components of $X$.
Note that for every purely inseparable algebraic extension $k'/k$
the morphism $(X_{k'})_{red} \to X$ is a homeomorphism, see
Algebra, Lemma \ref{algebra-lemma-p-ring-map}. Hence the points
$x'_1, \ldots, x'_n$ lying over $x_1, \ldots, x_n$ are the generic
points of the irreducible components of $(X_{k'})_{red}$.
As $X$ is reduced the local rings $K_i = \mathcal{O}_{X, x_i}$ are fields, see
Algebra, Lemma \ref{algebra-lemma-minimal-prime-reduced-ring}.
As $X$ is locally of finite type over $k$ the field extensions
$K_i/k$ are finitely generated field extensions.
Finally, the local rings $\mathcal{O}_{(X_{k'})_{red}, x'_i}$ are the
fields $(K_i \otimes_k k')_{red}$. By
Algebra, Lemma \ref{algebra-lemma-make-separable}
we can find a finite purely inseparable extension $k'/k$
such that $(K_i \otimes_k k')_{red}$ are separable field
extensions of $k'$. In particular each $(K_i \otimes_k k')_{red}$
is geometrically reduced over $k'$ by
Algebra, Lemma \ref{algebra-lemma-characterize-separable-field-extensions}.
At this point
Lemma \ref{lemma-generic-points-geometrically-reduced} part (3)
implies that $(X_{k'})_{red}$ is geometrically reduced.
\end{proof}

\end{document}
